% -------------
% Módulo con la configuración del documento y personalización de datos
% -------------

% Documento
% -------------

% Estilo de portada (0 - completa, 1 - cuarto de página, otro caso => 0)
\newcommand{\portada}{0}

% Tipografía para portada y por defecto (@Ver src/preambulo/fuentes)
\SetTypos[firasans][roboto]

% Paleta de colores (@Ver src/preambulo/temas)
\newcommand{\tema}{darkcula}

% Tema por defecto a \code (@See src/preambulo/style)
\lstset{style=mucho}


% Datos para portada
% -------------

% Nombre completo de la Universidad
\newcommand{\Universidad}{Nombre de la Universidad}

% Nombre de la Facultad
\newcommand{\Facultad}{Nombre de la Facultad}

% Nombre de la Carrera
\newcommand{\Carrera}{Nombre de la Carrera}

% Semestre actual
\newcommand{\Semestre}{2027-1}

% Nombre de la Materia
\newcommand{\Materia}{Nombre de la materia}

% Título completo de la obra
\newcommand{\Titulo}{Titulo Completo}

% Logos
%  Preferentemente agrega pdf
% --------------
% Nombre del logo 1 para DARK THEMES, preferentemente de tu universidad. Debe de estar en src/assets/logos_portada
\newcommand{\LogoUniD}{src/assets/logos_portada/uni_d.pdf}
% Nombre del logo 1 para LIGHT THEMES, preferentemente de tu universidad. Debe de estar en src/assets/logos_portada
\newcommand{\LogoUniL}{src/assets/logos_portada/uni_l.pdf}
% Nombre del logo 2 para DARK THEMES, preferentemente de tu facultad, intituto o carrera. Debe de estar en src/assets/logos_portada
\newcommand{\LogoFacD}{src/assets/logos_portada/fac_d.pdf}
% Nombre del logo 2 para LIGHT THEMES, preferentemente de tu facultad, intituto o carrera. Debe de estar en src/assets/logos_portada
\newcommand{\LogoFacL}{src/assets/logos_portada/fac_l.pdf}


% Nombre de los integrantes e identificador
\newcommand{\Integrantes}{
    Nombre Completo 1 & ID 1 \\
    Nombre Completo 2 & ID 2 \\
    Nombre Completo 3 & ID 3 \\
    \hline
}

% Nombre del profesor y puesto
\newcommand{\Docentes}{
    Nombre Completo del Profesor & Profesor \\
    Nombre Completo del Asistente & Asistente \\
    \hline
}

% Datos para header/footer (personaliza el tuyo en src/portadas/HeaderFooter.tex)
% -------------
% Siglas / Acrónimo de la Universidad (para header)
\newcommand{\UniversidadSiglas}{UNI}

% Siglas / Acrónimo de la Facultad (para header)
\newcommand{\FacultadSiglas}{FAC}

% Siglas / Acrónimo de la Carrera (para header)
\newcommand{\CarreraSiglas}{CARRERA}

% Siglas / Acrónimo de la materia (para header, si no tiene procurar una cadena corta)
\newcommand{\MateriaSiglas}{Materia}

% Versión acortada (si es necesario) del título (para header)
\newcommand{\TituloCompacto}{Titulo}

% Nombre de equipo o Apellidos de integrantes
\newcommand{\NombreCompacto}{Alumno 1. Alumno 2.}

% Copyrights
% -------------

% Título de la obra
\title{\Titulo}

% Nombre del autor
\author{\Integrantes}
